\chapter{تطبيقات نظرية الزمر}\label{ch:b3:group-theory-applied}

لا يشكّل ثنائي أكسيد الكربون إلا بضع مئات من الجزيئات في كل مليون من جزيئات
الهواء، والنيتروجين والأكسجين كل الباقي تقريبًا؛ ومع ذلك فإن ثنائي أكسيد الكربون، لا
النيتروجين ولا الأكسجين، هو الذي يمتص الأشعة تحت الحمراء التي ترسلها الأرض الدافئة نحو
الفضاء. فالجزيء لا يمتص الأشعة تحت الحمراء إلا عبر اهتزاز
يغيّر عزم ثنائي قطبه، وما إذا كان الاهتزاز يفعل ذلك يقرره
التناظر وحده: فاهتزاز النيتروجين الوحيد لا يستطيع، واثنان من اهتزازات ثنائي أكسيد الكربون
الأربعة يستطيعان. يحوّل هذا الفصل \omterm{def:b3:point-groups:irrep}{جداول المميِّزات} في
\cref{ch:b3:point-groups} إلى أدوات: فهو يختزل التمثيلات، ويبني
مدارات متكيفة مع التناظر، ويعدّ الاهتزازات ويسمّيها، ويستنتج
قواعد الانتقاء في مطيافية تحت الأحمر ومطيافية رامان.

\begin{recall}
عرّف \cref{ch:b3:point-groups} \omterm{def:b3:point-groups:operation}{عمليات التناظر} والزمر النقطية
والأصناف والتمثيلات \omterm{def:b3:point-groups:representation}{والمميِّزات} \omterm{def:b3:point-groups:irrep}{وجداول المميِّزات}. وبنى مجلد السنة الثانية
المدارات الجزيئية للجزيئات \ce{H2O} و\ce{NH3} و\ce{CH4} من
مدارات الشظايا ومن التركيبات المتكيفة مع التناظر لمدارات $1s$ للهيدروجين،
وسمّاها برموز تُستنتج هنا. وقرأ مجلد السنة الأولى
أطياف تحت الأحمر بالأعداد الموجية وعرّف قابلية استقطاب
الجزيء.
\end{recall}

\section{اختزال التمثيل}

كل مجموعة من الدوال أو المتجهات المرتبطة بجزيء تحمل
تمثيلًا \omterm{def:b3:point-groups:point-group}{لزمرته النقطية}. ومعظمها قابل للاختزال.

\begin{definition}[التمثيل القابل للاختزال، التمثيل المتناظر كليًا]\label{def:b3:group-theory-applied:reducible}
\emph{التمثيل القابل للاختزال}\index{تمثيل قابل للاختزال} تمثيل
يمكن أن تُردّ مصفوفاته كلها، بتغيير واحد للأساس، إلى الشكل
القطري بالكتل نفسه؛ وهو عندئذ مجموع \omterm{def:b3:point-groups:irrep}{تمثيلات غير قابلة للاختزال}،
يُكتب $\Gamma = \sum_in_i\Gamma_i$. \omterm{def:b3:point-groups:irrep}{والتمثيل غير القابل للاختزال} الذي
\omterm{def:b3:point-groups:representation}{مميِّزاته} كلها 1 هو \emph{التمثيل المتناظر
كليًا}\index{تمثيل متناظر كليًا} ($A_1$، $A_{1g}$،
$A_1'$\dots).
\end{definition}

\begin{theorem}[مبرهنة التعامد الكبرى]\label{thm:b3:group-theory-applied:great-orthogonality}
من أجل تمثيلين غير قابلين للاختزال $\Gamma_i$ و$\Gamma_j$ بعداهما $d_i$
و$d_j$ لزمرة رتبتها $h$، مكتوبين بمصفوفات واحدية،
\[
\sum_R D_i(R)_{mn}^*\,D_j(R)_{m'n'} = \frac{h}{d_i}\,\delta_{ij}\,\delta_{mm'}\,\delta_{nn'}.
\]
\end{theorem}
\admitted

البرهان، بتوطئتي شور، يُتناول في مقررات أكثر تقدمًا؛ أما
نتائجه في \omterm{def:b3:point-groups:representation}{المميِّزات} فهي ما تستعمله الكيمياء، وكل جدول \omterm{def:b3:point-groups:representation}{مميِّزات}
في هذا الكتاب مُتحقَّق منه بها.

\begin{corollary}[تعامد المميِّزات]\label{cor:b3:group-theory-applied:characters}
إذا كان $g_c$ عدد العمليات في الصنف $c$، فإن
\[
\sum_cg_c\,\chi_i(c)^*\chi_j(c) = h\,\delta_{ij},
\]
وعدد \omterm{def:b3:point-groups:irrep}{التمثيلات غير القابلة للاختزال} يساوي عدد الأصناف،
مع $\sum_id_i^2 = h$.
\end{corollary}

\begin{proof}
لنضع $m = n$ و$m' = n'$ في المبرهنة ولنجمع على $m$ و$m'$: فيصير الطرف
الأيسر $\sum_R\chi_i(R)^*\chi_j(R)$، والطرف الأيمن $\frac{h}{d_i}\delta_{ij}
\sum_{m}\delta_{mm} = h\delta_{ij}$؛ ثم نجمّع العمليات بحسب الصنف. فأسطر
الجدول إذن متجهات متعامدة في فضاء بعده عدد
الأصناف، فلا تزيد \omterm{def:b3:point-groups:irrep}{التمثيلات غير القابلة للاختزال} على عدد
الأصناف؛ أما المساواة و$\sum d_i^2 = h$ فتأتيان من تطبيق
المبرهنة على التمثيل المنتظم (التمرين 10).
\end{proof}

\begin{theorem}[صيغة الاختزال]\label{thm:b3:group-theory-applied:reduction}
التمثيل $\Gamma$ ذو \omterm{def:b3:point-groups:representation}{المميِّزات} $\chi(c)$ يحتوي \omterm{def:b3:point-groups:irrep}{التمثيل غير القابل للاختزال}
$\Gamma_i$ بالضبط
\[
n_i = \frac1h\sum_cg_c\,\chi(c)\,\chi_i(c)^*
\]
مرة. هذه هي \emph{صيغة الاختزال}\index{صيغة الاختزال}.
\end{theorem}

\begin{proof}
بما أن $\Gamma = \sum_jn_j\Gamma_j$، \omterm{def:b3:point-groups:representation}{فمميِّزاته} $\chi(c) = \sum_jn_j\chi_j(c)$.
لنضرب في $g_c\chi_i(c)^*$، ولنجمع على الأصناف، ولنستعمل النتيجة:
$\sum_cg_c\chi(c)\chi_i(c)^* = \sum_jn_jh\delta_{ij} = hn_i$.
\end{proof}

\begin{example}[مدارات الهيدروجين في الأمونيا]\label{ex:b3:group-theory-applied:nh3}
لنأخذ مدارات $1s$ الثلاثة لذرات الهيدروجين في \ce{NH3} أساسًا. تترك $E$
الثلاثة في مواضعها ($\chi = 3$)، وتحرّك $C_3$ الثلاثة ($\chi = 0$)، ويترك
$\sigma_v$ واحدًا في موضعه ويبادل اثنين ($\chi = 1$). وبجدول $C_{3v}$:
$n_{A_1} = \frac16(3 + 0 + 3\times1) = 1$، $n_{A_2} = \frac16(3 + 0 - 3) = 0$،
$n_E = \frac16(6 + 0 + 0) = 1$. ومنه $\Gamma_{\mathrm H} = A_1 + E$: تعطي ذرات الهيدروجين
الثلاث تركيبًا واحدًا متناظرًا كليًا وزوجًا منحلًا.
\end{example}

\begin{method}[اختزال تمثيل]\label{met:b3:group-theory-applied:reduce}
\begin{enumerate}
  \item جد \omterm{def:b3:point-groups:representation}{المميِّز} لكل صنف: من أجل مجموعة من الدوال المتمركزة على
        الذرات، هو عدد الدوال التي تبقى في مواضعها، كل منها معدودة
        بالإشارة التي تكتسبها.
  \item طبّق \omterm{thm:b3:group-theory-applied:reduction}{صيغة الاختزال} بسطر كل \omterm{def:b3:point-groups:irrep}{تمثيل غير قابل
        للاختزال}، مرجِّحًا كل صنف بحجمه.
  \item تحقق: النتائج أعداد صحيحة غير سالبة، و$\sum_in_id_i$
        يساوي \omterm{def:b3:point-groups:representation}{المميِّز} تحت $E$ (حجم الأساس).
\end{enumerate}
\end{method}

\section{المدارات المتكيفة مع التناظر}

\begin{definition}[مؤثر الإسقاط]\label{def:b3:group-theory-applied:projection}
\emph{مؤثر الإسقاط}\index{مؤثر الإسقاط} على \omterm{def:b3:point-groups:irrep}{التمثيل
غير القابل للاختزال} $\Gamma_i$ هو
\[
\hat P_i = \frac{d_i}{h}\sum_R\chi_i(R)^*\,\hat R,
\]
حيث يطبّق $\hat R$ العملية $R$ على دالة.
\end{definition}

\begin{proposition}[ما يفعله مؤثر الإسقاط]\label{prop:b3:group-theory-applied:projection}
إذا طُبّق على أي دالة $f$، كان $\hat P_if$ إما معدومًا وإما دالة
تتحول كما يتحول $\Gamma_i$ (تركيبًا متكيفًا مع التناظر)؛ ومن أجل $\Gamma_i$
أحادي البعد لا تغيّره، بإهمال الإشارة $\chi_i(S)$، أي
عملية $S$.
\end{proposition}

\begin{proof}[برهان جزئي]
من أجل $d_i = 1$: $\hat S\hat P_if = \frac1h\sum_R\chi_i(R)\hat S\hat Rf$. لنضع $R' = SR$،
الذي يمسح الزمرة كلها كما يمسحها $R$؛ و$\chi_i(R) = \chi_i(S^{-1}R') =
\chi_i(S)\chi_i(R')$ من أجل تمثيل أحادي البعد (\omterm{def:b3:point-groups:representation}{فالمميِّزات} هي
المصفوفات نفسها). ومنه $\hat S\hat P_if = \chi_i(S)\hat P_if$. أما الحالة
العامة فتستعمل مبرهنة التعامد الكبرى ونقبلها دون برهان.
\end{proof}

\begin{method}[بناء التركيبات المتكيفة مع التناظر]\label{met:b3:group-theory-applied:salc}
\begin{enumerate}
  \item اختر دالة واحدة من المجموعة، $h_1$، واكتب في جدول الدالة التي ترسلها
        إليها كل عملية.
  \item من أجل كل \omterm{def:b3:point-groups:irrep}{تمثيل غير قابل للاختزال}، اضرب كل صورة في
        \omterm{def:b3:point-groups:representation}{مميِّز} العملية، واجمع، ثم نظّم (مهملًا التداخل
        بين الدوال الذرية).
  \item من أجل تمثيل منحل، أسقِط دالة ثانية للحصول على
        عضو ثانٍ، ثم عامِد.
\end{enumerate}
\end{method}

\begin{example}[تركيبات الأمونيا والماء]\label{ex:b3:group-theory-applied:salcs}
في \ce{NH3}، $\hat P_{A_1}h_1 \propto h_1 + h_2 + h_3$ (كل \omterm{def:b3:point-groups:representation}{المميِّزات} 1). ومن أجل $E$
(\omterm{def:b3:point-groups:representation}{المميِّزات} $2, -1, 0$)، $\hat P_Eh_1 \propto 2h_1 - h_2 - h_3$؛ وإسقاط $h_2$ ثم
المعامدة يعطيان الشريك $h_2 - h_3$. وبعد التنظيم: $a_1 =
\frac{1}{\sqrt3}(h_1 + h_2 + h_3)$، $e = \frac{1}{\sqrt6}(2h_1 - h_2 - h_3)$،
$\frac{1}{\sqrt2}(h_2 - h_3)$. وفي \ce{H2O}، الموضوع في المستوي $xz$، يكون $h_1 + h_2$
من النوع $a_1$ و$h_1 - h_2$ من النوع $b_1$ (فهو يغيّر إشارته بالعمليتين $C_2$ و$\sigma_v'(yz)$،
اللتين تبادلان ذرتي الهيدروجين). وإذا وُضع الجزيء في المستوي $yz$، كما تفضّل بعض
الكتب، تبادل الرمزان $b_1$ و$b_2$ في كل مكان.
\end{example}

\begin{omfigure}
\begin{tikzpicture}[font=\small]
  \foreach \x/\lab/\ca/\cb/\cc in {0/{$a_1$}/1/1/1, 4/{$e$}/2/-1/-1, 8/{$e$}/0/1/-1}{
    \begin{scope}[xshift=\x cm]
      \foreach \a/\c in {90/\ca,210/\cb,330/\cc}{
        \pgfmathsetmacro{\r}{0.18+0.13*abs(\c)}
        \ifdim \c pt > 0pt \fill[omDef!55, draw=omDef] (\a:1.0) circle (\r);
        \else \ifdim \c pt < 0pt \fill[omThm!25, draw=omThm] (\a:1.0) circle (\r);
        \else \draw[gray, densely dotted] (\a:1.0) circle (0.12); \fi\fi
      }
      \node[aN, minimum size=4.6mm] at (0,0) {};
      \node[anchor=north] at (0,-1.45) {\lab};
    \end{scope}}
  \node[anchor=north, font=\footnotesize] at (0,-1.85) {$h_1 + h_2 + h_3$};
  \node[anchor=north, font=\footnotesize] at (4,-1.85) {$2h_1 - h_2 - h_3$};
  \node[anchor=north, font=\footnotesize] at (8,-1.85) {$h_2 - h_3$};
\end{tikzpicture}

{\small التركيبات المتكيفة مع التناظر لمدارات $1s$ الثلاثة للهيدروجين
في الأمونيا، منظورًا إليها على امتداد المحور $C_3$ ($h_1$ في الأعلى). الأزرق: معامل
موجب، والأحمر: سالب، وحجم كل قرص يكبر مع
المعامل؛ والمنقّط: صفر.}
\end{omfigure}

\begin{example}[ست ربائط حول فلز]\label{ex:b3:group-theory-applied:ml6}
من أجل ستة مدارات $\sigma$ للربائط متجهة نحو فلز من رؤوس
ثماني أوجه، تكون \omterm{def:b3:point-groups:representation}{المميِّزات} على أصناف $O_h$ ($E$، $8C_3$، $6C_2$، $6C_4$،
$3C_2$، $i$، $6S_4$، $8S_6$، $3\sigma_h$، $6\sigma_d$) هي $6, 0, 0, 2, 2, 0, 0, 0, 4,
2$، ويعطي الاختزال $\Gamma_\sigma = A_{1g} + E_g + T_{1u}$. فتجد مدارات الفلز $s$
($a_{1g}$) و$d_{z^2}$ و$d_{x^2-y^2}$ ($e_g$) و$p$ ($t_{1u}$) شركاء من الربائط؛
أما مداراته $d_{xy}$ و$d_{xz}$ و$d_{yz}$ ($t_{2g}$) فلا تجد شريكًا وتبقى غير رابطة
بالنسبة إلى $\sigma$ --- وهذا منشأ الانشطار $t_{2g}$/$e_g$ في مجلد السنة الثانية،
الذي يُطوَّر في \cref{ch:b3:organometallic-bonding,ch:b3:complex-spectra-magnetism}.
\end{example}

\begin{omfigure}
\begin{tikzpicture}[font=\small, yscale=0.95]
  % O orbitals (left), MOs (centre), H combinations (right); schematic energies
  \foreach \y/\l in {0/{$2s\ (a_1)$}}{\draw[omstate] (0,\y) -- (1.2,\y); \node[anchor=east] at (-0.05,\y) {\l};}
  \draw[omstate] (0,3.0) -- (0.36,3.0) (0.42,3.0) -- (0.78,3.0) (0.84,3.0) -- (1.2,3.0);
  \node[anchor=east] at (-0.05,3.0) {$2p\ (b_1, a_1, b_2)$};
  \draw[omstate] (6.8,3.9) -- (8.0,3.9); \node[anchor=west] at (8.05,3.9) {$h_1 - h_2\ (b_1)$};
  \draw[omstate] (6.8,3.6) -- (8.0,3.6); \node[anchor=west] at (8.05,3.5) {$h_1 + h_2\ (a_1)$};
  \foreach \y/\l in {-0.6/{$2a_1$}, 1.6/{$1b_1$}, 2.4/{$3a_1$}, 3.0/{$1b_2$}, 5.6/{$4a_1$}, 6.2/{$2b_1$}}{
    \draw[omstate] (3.4,\y) -- (4.6,\y); \node[anchor=west] at (4.65,\y) {\l};}
  \foreach \y in {-0.6,1.6,2.4,3.0}{\node[font=\scriptsize] at (4.0,\y+0.16) {$\uparrow\downarrow$};}
  \draw[densely dotted, gray] (1.2,0) -- (3.4,-0.6) (1.2,3.0) -- (3.4,1.6) (1.2,3.0) -- (3.4,2.4) (1.2,3.0) -- (3.4,3.0)
     (6.8,3.6) -- (4.6,-0.6) (6.8,3.9) -- (4.6,1.6) (6.8,3.6) -- (4.6,2.4) (1.2,3.0) -- (3.4,6.2) (6.8,3.9) -- (4.6,6.2) (6.8,3.6) -- (4.6,5.6);
  \node[anchor=north] at (0.6,-1.0) {O};
  \node[anchor=north] at (4.0,-1.0) {\ce{H2O}};
  \node[anchor=north] at (7.4,-1.0) {2 H};
  \draw[->] (-2.3,-0.9) -- (-2.3,6.4) node[above] {$E$};
\end{tikzpicture}

{\small مدارات التكافؤ الجزيئية للماء (طاقات تخطيطية، والجزيء في
المستوي $xz$). لا تمتزج إلا المدارات ذات التناظر نفسه: تركيب الهيدروجين $b_1$
مع $2p_x$، وتركيب $a_1$ مع $2s$ و$2p_z$؛ أما $2p_y$ ($b_2$) فلا
شريك له ويبقى زوجًا حرًا غير رابط، وهو أعلى مدار مشغول.}
\end{omfigure}

\section{أنماط الاهتزاز}

\begin{definition}[النمط العادي]\label{def:b3:group-theory-applied:normal-mode}
\emph{النمط العادي}\index{نمط عادي} للجزيء اهتزاز
جماعي تهتز فيه كل الذرات بالتردد نفسه وبالطور نفسه،
على امتداد المتجهة الذاتية لمصفوفة هيس المرجَّحة بالكتلة
(\cref{prop:b3:computational-chemistry:hessian}). وكل نمط عادي
يتحول كما يتحول \omterm{def:b3:point-groups:irrep}{تمثيل غير قابل للاختزال} \omterm{def:b3:point-groups:point-group}{للزمرة النقطية}.
\end{definition}

\begin{proposition}[تمثيل الحركات كلها]\label{prop:b3:group-theory-applied:gamma-3n}
لنُلحق بكل ذرة ثلاث متجهات إزاحة $x$ و$y$ و$z$. تحت عملية
$R$ يكون \omterm{def:b3:point-groups:representation}{مميِّز} هذا التمثيل ذي البعد $3N$
\[
\chi_{3N}(R) = (\text{عدد الذرات الباقية في مواضعها}) \times \chi_{xyz}(R),
\]
حيث $\chi_{xyz} = 1 + 2\cos\theta$ لدوران بزاوية $\theta$ و$-1 + 2\cos\theta$
لدوران غير فعلي (أي 3 للعملية $E$، و$-1$ للعملية $C_2$، و0 للعملية $C_3$، و1 للعملية $\sigma$،
و$-3$ للعملية $i$).
\end{proposition}

\begin{proof}
الذرة المنقولة إلى موضع آخر تحمل متجهاتها خارج قطر
المصفوفة: فلا تساهم بشيء في الأثر. والذرة الباقية في موضعها
تتحول متجهاتها الثلاث بالمصفوفة $3\times3$ للعملية $R$، التي يُحسب أثرها
في معلم يكون فيه $z$ على امتداد المحور: $\cos\theta + \cos\theta + 1$ من أجل
دوران، ويصير العنصر الأخير $-1$ من أجل دوران غير فعلي.
\end{proof}

\begin{proposition}[عدّ الاهتزازات]\label{prop:b3:group-theory-applied:mode-count}
$\Gamma_{\mathrm{vib}} = \Gamma_{3N} - \Gamma_{\mathrm{trans}} - \Gamma_{\mathrm{rot}}$، حيث
$\Gamma_{\mathrm{trans}}$ تمثيل $(x, y, z)$ و$\Gamma_{\mathrm{rot}}$ تمثيل
$(R_x, R_y, R_z)$؛ وللجزيء غير الخطي $3N - 6$ اهتزازًا، وللجزيء
الخطي $3N - 5$.
\end{proposition}

\begin{proof}
تصف الإزاحات $3N$ كل حركة: ثلاثة انسحابات لمركز
الكتلة، وثلاثة دورانات (اثنان للجزيء الخطي، لأن دورانه
حول محوره لا يحرّك أي نواة)، والاهتزازات. وهذه الفضاءات الجزئية لا
تمتزج تحت العمليات، فتتجمع تمثيلاتها بالجمع.
\end{proof}

\begin{example}[الماء]\label{ex:b3:group-theory-applied:water}
% ledger: vib:H2O.sym, vib:H2O.bend, vib:H2O.asym
\ce{H2O} في $C_{2v}$، والجزيء في المستوي $xz$. الذرات الباقية في مواضعها: 3، 1، 3، 1؛
و$\chi_{xyz}$: $3, -1, 1, 1$؛ ومنه $\chi_{3N} = 9, -1, 3, 1$. ويعطي الاختزال $3A_1 +
A_2 + 3B_1 + 2B_2$. وبطرح الانسحابات $A_1 + B_1 + B_2$ والدورانات $A_2 + B_1 + B_2$
يبقى $\Gamma_{\mathrm{vib}} = 2A_1 + B_1$. النمطان $a_1$ هما التمدد
المتناظر (\qty{3657}{cm^{-1}}) والانحناء (\qty{1595}{cm^{-1}})؛ والنمط $b_1$ هو
التمدد غير المتناظر (\qty{3756}{cm^{-1}}).
\end{example}

\begin{omfigure}
\begin{tikzpicture}[font=\small, >=Stealth]
  \foreach \x/\name/\lab in {0/sym/{\foreignlanguage{arabic}{$\nu_1$، $a_1$، تمدد متناظر}}, 4.3/bend/{\foreignlanguage{arabic}{$\nu_2$، $a_1$، انحناء}}, 8.6/asym/{\foreignlanguage{arabic}{$\nu_3$، $b_1$، تمدد غير متناظر}}}{
    \begin{scope}[xshift=\x cm]
      \node[aO] (o\name) at (0,0.35) {};
      \node[aH] (a\name) at (-0.8,-0.25) {};
      \node[aH] (b\name) at (0.8,-0.25) {};
      \begin{scope}[on background layer]\draw[bond] (o\name.center) -- (a\name.center) (o\name.center) -- (b\name.center);\end{scope}
      \node[anchor=north, align=center, font=\footnotesize] at (0,-0.9) {\lab};
    \end{scope}}
  % symmetric stretch: H outward along the bonds, O up a little
  \draw[->, omThm, thick] (-0.95,-0.36) -- (-1.42,-0.71);
  \draw[->, omThm, thick] (0.95,-0.36) -- (1.42,-0.71);
  \draw[->, omThm, thick] (0,0.74) -- (0,0.98);
  % bend: each H moves perpendicular to its bond, closing the angle; O up
  \draw[->, omThm, thick] (4.3-0.92,-0.36) -- (4.3-0.62,-0.76);
  \draw[->, omThm, thick] (4.3+0.92,-0.36) -- (4.3+0.62,-0.76);
  \draw[->, omThm, thick] (4.3,0.74) -- (4.3,0.98);
  % antisymmetric stretch: one H out, the other in, O sideways
  \draw[->, omThm, thick] (8.6-0.95,-0.36) -- (8.6-1.42,-0.71);
  \draw[->, omThm, thick] (8.6+1.02,0.0) -- (8.6+0.70,0.24);
  \draw[->, omThm, thick] (8.6-0.05,0.80) -- (8.6+0.25,0.80);
\end{tikzpicture}

{\small \omterm{def:b3:group-theory-applied:normal-mode}{الأنماط العادية} الثلاثة للماء (الأسهم: إزاحات، ليست
بمقياس الرسم). يحفظ النمطان $a_1$ كلاهما التناظر الكامل للجزيء؛ أما النمط
$b_1$ فيغيّر إشارته بالعملية $C_2$.}
\end{omfigure}

\begin{example}[أربعة جزيئات أخرى]\label{ex:b3:group-theory-applied:table}
يعطي الإجراء نفسه (المحسوب والمتحقَّق منه في معطيات أشكال هذا
الفصل):
\begin{center}\small
\begin{tabular}{llll}
\hline
الجزيء & الزمرة & $\Gamma_{\mathrm{vib}}$ & الأنماط \\ \hline
\ce{NH3} & $C_{3v}$ & $2A_1 + 2E$ & 6 \\
\ce{CH4} & $T_d$ & $A_1 + E + 2T_2$ & 9 \\
\ce{CO2} & $D_{\infty h}$ (عبر $D_{2h}$) & $\Sigma_g^+ + \Sigma_u^+ + \Pi_u$ ($A_g + B_{1u} + B_{2u} + B_{3u}$) & 4 \\
\ce{BF3} & $D_{3h}$ & $A_1' + 2E' + A_2''$ & 6 \\
\ce{XeF4} & $D_{4h}$ & $A_{1g} + B_{1g} + B_{2g} + A_{2u} + B_{2u} + 2E_u$ & 9 \\
\ce{SF6} & $O_h$ & $A_{1g} + E_g + T_{2g} + 2T_{1u} + T_{2u}$ & 15 \\ \hline
\end{tabular}
\end{center}
من أجل جزيء خطي تُستبدل بالزمرة اللانهائية زمرتها الجزئية $D_{2h}$،
التي ينشطر فيها الانحناء المنحل $\Pi_u$ إلى $B_{2u} + B_{3u}$.
\end{example}

\section{قواعد الانتقاء}

\begin{definition}[الجداء المباشر]\label{def:b3:group-theory-applied:direct-product}
\emph{الجداء المباشر}\index{جداء مباشر} $\Gamma_i\otimes\Gamma_j$
لتمثيلين هو التمثيل الذي تحمله جداءات دوالهما
الأساسية؛ \omterm{def:b3:point-groups:representation}{ومميِّزاته} هي الجداءات $\chi_i(R)\chi_j(R)$.
\end{definition}

\begin{theorem}[التكاملات المعدومة]\label{thm:b3:group-theory-applied:vanishing-integral}
لا يمكن أن يكون التكامل $\int f_1f_2f_3\,\dd\tau$ على الفضاء كله غير معدوم إلا إذا
احتوى \omterm{def:b3:group-theory-applied:direct-product}{الجداء المباشر} $\Gamma_1\otimes\Gamma_2\otimes\Gamma_3$ \omterm{def:b3:group-theory-applied:reducible}{التمثيل المتناظر
كليًا}.
\end{theorem}

\begin{proof}
لا تفعل \omterm{def:b3:point-groups:operation}{عملية التناظر} إلا إعادة تسمية نقاط الفضاء، فلا تتغير قيمة
التكامل حين تُطبَّق أي عملية $R$ على المكامَل. لنأخذ متوسط
المكامَل المحوَّل على الزمرة: متوسط مركبات
المكامَل المنتمية إلى كل \omterm{def:b3:point-groups:irrep}{تمثيل غير قابل للاختزال} هو
$\frac1h\sum_R\hat R(\cdot)$، أي الإسقاط على \omterm{def:b3:group-theory-applied:reducible}{التمثيل المتناظر
كليًا} (\cref{def:b3:group-theory-applied:projection} مع كل
$\chi = 1$). وكل مركبة أخرى متوسطها صفر؛ فإذا لم يحتو الجداء أي
جزء متناظر كليًا، كان التكامل معدومًا.
\end{proof}

\begin{definition}[النشاط في تحت الأحمر، النشاط في رامان]\label{def:b3:group-theory-applied:activity}
يكون الاهتزاز الأساسي \emph{نشطًا في تحت الأحمر}\index{نشط في تحت الأحمر} إذا استطاع أن يمتص
الأشعة تحت الحمراء، و\emph{نشطًا في رامان}\index{نشط في رامان} إذا ظهر في
طيف رامان (\cref{ch:b3:rovibrational-spectroscopy}).
\end{definition}

\begin{corollary}[قواعد الانتقاء في تحت الأحمر ورامان]\label{cor:b3:group-theory-applied:ir-raman}
لا يكون الانتقال الأساسي (من المستوى الأساسي، المتناظر كليًا، إلى كمّ واحد
من نمط تناظره $\Gamma$) \omterm{def:b3:group-theory-applied:activity}{نشطًا في تحت الأحمر} إلا إذا كان $\Gamma$
تمثيل $x$ أو $y$ أو $z$، ولا \omterm{def:b3:group-theory-applied:activity}{نشطًا في رامان} إلا إذا كان $\Gamma$
تمثيل دالة تربيعية ($x^2$، $xy$، \dots).
\end{corollary}

\begin{proof}
تتضمن شدة الامتصاص التكامل $\int\psi_0\,\mu_k\,\psi_1\dd\tau$ مع مركبات ثنائي
القطب $\mu_k$، التي تتحول كما تتحول $x$ و$y$ و$z$؛ و$\psi_0$ متناظرة كليًا
و$\psi_1$ تتحول كما يتحول $\Gamma$. وبحسب المبرهنة لا يمكن أن يكون التكامل غير معدوم
إلا إذا احتوى $\Gamma\otimes\Gamma_{\mu_k}$ على $A_1$، وهذا لا يحدث إلا إذا كان $\Gamma =
\Gamma_{\mu_k}$ (فجداء تمثيلين غير قابلين للاختزال لا يحتوي \omterm{def:b3:group-theory-applied:reducible}{التمثيل
المتناظر كليًا} إلا إذا تساويا، من أجل \omterm{def:b3:point-groups:representation}{مميِّزات} حقيقية). أما انتثار
رامان فيتضمن مركبات قابلية الاستقطاب $\alpha_{kl}$، التي
تتحول كما تتحول الجداءات $kl$.
\end{proof}

\begin{proposition}[قاعدة الاستبعاد المتبادل]\label{prop:b3:group-theory-applied:mutual-exclusion}
في جزيء له \omterm{def:b3:point-groups:inversion}{مركز انقلاب}، لا يكون أي انتقال أساسي \omterm{def:b3:group-theory-applied:activity}{نشطًا في تحت الأحمر} ورامان
معًا: هذه \emph{قاعدة الاستبعاد المتبادل}\index{قاعدة الاستبعاد المتبادل}.
\end{proposition}

\begin{proof}
مع وجود \omterm{def:b3:point-groups:inversion}{مركز انقلاب} يكون كل \omterm{def:b3:point-groups:irrep}{تمثيل غير قابل للاختزال} من النوع $g$ أو $u$. 
والإحداثيات $x$ و$y$ و$z$ تغيّر إشارتها تحت $i$ (فهي $u$)؛ أما الدوال التربيعية
فلا ($g$). فلا يكون النمط \omterm{def:b3:group-theory-applied:activity}{نشطًا في تحت الأحمر} إلا إذا كان $u$، ولا \omterm{def:b3:group-theory-applied:activity}{نشطًا في رامان} إلا إذا كان
$g$.
\end{proof}

\begin{example}[ثنائي أكسيد الكربون والدفيئة]\label{ex:b3:group-theory-applied:co2}
% ledger: vib:CO2.sym, vib:CO2.bend, vib:CO2.asym, ir:CO2.asym.int, ir:CO2.bend.int
لثنائي أكسيد الكربون \ce{CO2} \omterm{def:b3:point-groups:inversion}{مركز انقلاب}. فتمدده المتناظر ($\Sigma_g^+$،
\qty{1333}{cm^{-1}}) \omterm{def:b3:group-theory-applied:activity}{نشط في رامان} وحده؛ أما التمدد غير المتناظر ($\Sigma_u^+$،
\qty{2349}{cm^{-1}}) والانحناء ($\Pi_u$، \qty{667}{cm^{-1}}) فنشطان في تحت الأحمر وحده.
ويمتص الانحناء قرب \qty{15}{\micro m}، قريبًا من ذروة الإصدار الحراري
للأرض: وهذا سبب كون ثنائي أكسيد الكربون غاز دفيئة. أما \ce{N2} 
و\ce{O2} فلكل منهما اهتزاز واحد، من النوع $\Sigma_g^+$: غير \omterm{def:b3:group-theory-applied:activity}{نشط في تحت الأحمر}.
\end{example}

\begin{omfigure}
\begin{tikzpicture}
\begin{axis}[omir, width=0.92\linewidth, height=4.4cm, xmin=400, xmax=2600,
  ymin=0, ymax=1.3, ytick=\empty, xlabel={\foreignlanguage{arabic}{العدد الموجي \babelsublr{(cm$^{-1}$)}}},
  ylabel={\foreignlanguage{arabic}{الإشارة}}, legend cell align=left,
  legend style={font=\small, at={(0.98,0.95)}, anchor=north east, draw=none}]
\addplot[omDef, thick] table[x=nu, y=ir] {figdata/out/bachelor-3-co2-spectra.dat};
\addlegendentry{\foreignlanguage{arabic}{امتصاص تحت الأحمر}}
\addplot[omThm, thick, dashed] table[x=nu, y=raman] {figdata/out/bachelor-3-co2-spectra.dat};
\addlegendentry{\foreignlanguage{arabic}{انتثار رامان}}
\node[font=\small, anchor=south] at (axis cs:2349,1.02) {$\Sigma_u^+$};
\node[font=\small, anchor=south] at (axis cs:667,0.1) {$\Pi_u$};
\node[font=\small, anchor=south, omThm] at (axis cs:1333,1.02) {$\Sigma_g^+$};
\end{axis}
\end{tikzpicture}

{\small أطياف اصطناعية لغاز \ce{CO2} (مواضع الحزم وشداتها النسبية
في تحت الأحمر مأخوذة من القيم المقيسة؛ والأشكال تخطيطية، دون
بنية دورانية). الاستبعاد المتبادل: لا تظهر أي حزمة في الطيفين معًا.}
\end{omfigure}

\begin{method}[التنبؤ بطيف اهتزازي]\label{met:b3:group-theory-applied:vibrations}
\begin{enumerate}
  \item جد \omterm{def:b3:point-groups:point-group}{الزمرة النقطية}؛ واحسب $\chi_{3N}$ من الذرات الباقية في مواضعها.
  \item اختزل، واطرح الانسحابات والدورانات: $\Gamma_{\mathrm{vib}}$.
  \item من الأعمدة اليمنى للجدول، علّم كل نمط بأنه \omterm{def:b3:group-theory-applied:activity}{نشط في تحت الأحمر}
        ($x$، $y$، $z$) و/أو \omterm{def:b3:group-theory-applied:activity}{نشط في رامان} (تربيعي)؛ والنمط الذي ليس هذا
        ولا ذاك صامت.
  \item عُدّ الحزم: حزمة واحدة لكل نمط نشط، والنمط المنحل يعطي حزمة واحدة.
\end{enumerate}
\end{method}

\begin{method}[عدّ تمددات CO]\label{met:b3:group-theory-applied:co-count}
في كربونيل فلزي، خذ متجهات الروابط C--O وعددها $n$ أساسًا (الذرة الباقية في
موضعها تُعدّ 1، ولا شيء غير ذلك): اختزل للحصول على $\Gamma_{\mathrm{CO}}$، ثم طبّق
قاعدتي تحت الأحمر ورامان. فعدد الحزم القوية قرب
\qtyrange{1850}{2150}{cm^{-1}} يحدد المتماكب.
\end{method}

\begin{example}[cis وtrans]\label{ex:b3:group-theory-applied:cis-trans}
cis-\ce{ML4(CO)2}، $C_{2v}$: \omterm{def:b3:point-groups:representation}{المميِّزات} $2, 0, 2, 0$، $\Gamma_{\mathrm{CO}} = A_1 + B_1$،
حزمتان في تحت الأحمر. trans-\ce{ML4(CO)2}، $D_{4h}$: $\Gamma_{\mathrm{CO}} = A_{1g} + A_{2u}$، حزمة
واحدة في تحت الأحمر ($A_{2u}$) وحزمة واحدة في رامان ($A_{1g}$). فحزمة CO واحدة في
تحت الأحمر تعني trans.
\end{example}

\section{التطبيقات على الانتقالات الإلكترونية}

تنطبق المبرهنة نفسها على الانتقالات الإلكترونية: فالانتقال الإلكتروني
بين الحالتين $\Psi_1$ و$\Psi_2$ لا يكون مسموحًا إلا إذا احتوى $\Gamma_1\otimes\Gamma_2$
تمثيل $x$ أو $y$ أو $z$، ويكون ضوؤه مستقطبًا على امتداد
ذلك المحور.

\begin{example}[انتقالات الماء والميثانال]\label{ex:b3:group-theory-applied:electronic}
في الماء ($C_{2v}$، المستوي $xz$)، يعطي ترقية إلكترون من $1b_2$ (الزوج الحر)
إلى $4a_1$ حالة مثارة تناظرها $B_2\otimes A_1 = B_2$، وهو $y$:
فالانتقال مسموح، ومستقطب عموديًا على مستوي الجزيء. وفي الميثانال
($C_{2v}$)، ينطلق الانتقال $n\to\pi^*$ من زوج حر للأكسجين واقع في المستوي ($b_1$)
إلى المدار $\pi^*$، المبني من مدارات $p$ عمودية على المستوي
($b_2$؛ والجزيء في المستوي $xz$، وC=O على امتداد $z$): $B_1\otimes B_2 = A_2$،
وهو ليس أيًا من $x$ و$y$ و$z$: فهو ممنوع بالتناظر، ولهذا تكون حزمته
(\cref{ch:b3:electronic-spectroscopy}) ضعيفة جدًا.
\end{example}

\begin{history}[فيغنر وبيته وتناظر الجزيئات]
دخلت نظرية الزمر ميكانيكا الكم بأعمال يوجين فيغنر في الأطياف
الذرية (1927--31) وبمقالة هانز بيته في انشطار الحدود الذرية في
البلورات (1929). وطبّقها إ. برايت ويلسون على الاهتزازات الجزيئية سنة
1934؛ وأعطى روبرت موليكن الرموز التي ما زالت تُستعمل للمدارات والحالات.
\end{history}

\begin{inthelab}[تحت الأحمر ورامان جنبًا إلى جنب]
يقيس مطياف تحت الأحمر الضوء الذي تمتصه العيّنة؛ أما مطياف
رامان فيضيئها بليزر ويحلل الضوء المنتثر الضعيف
عند أطوال موجية منزاحة. وتسجيل الطيفين كليهما للمركب نفسه هو
الاختبار الكلاسيكي لوجود مركز تناظر: فإذا لم تتطابق أي حزمة، فالجزيء
على الأرجح متناظر مركزيًا. وليزرات رامان من الصنف 3B أو 4: فتُحبس الحزمة
الضوئية ويضع العاملون نظارات واقية مناسبة لطول موجتها.
\end{inthelab}

\section{تمارين}

\begin{exercise}[$\star$]\label{exo:b3:group-theory-applied:1}
اختزل تمثيل $C_{2v}$ ذا \omterm{def:b3:point-groups:representation}{المميِّزات} $(5, -1, 3, 1)$ على $(E, C_2,
\sigma_v(xz), \sigma_v'(yz))$. وتحقق من البعد.
\end{exercise}

\begin{exercise}[$\star$]\label{exo:b3:group-theory-applied:2}
الجزيء \ce{SO2} منحنٍ، كالماء. أعطِ $\Gamma_{\mathrm{vib}}$ وبيّن أي الأنماط \omterm{def:b3:group-theory-applied:activity}{نشط في تحت الأحمر}
وأيها \omterm{def:b3:group-theory-applied:activity}{نشط في رامان}.
\end{exercise}

\begin{exercise}[$\star$]\label{exo:b3:group-theory-applied:3}
في $C_{2v}$، احسب الجداءين المباشرين $B_1\otimes B_2$ و$A_2\otimes B_1$، 
و$E\otimes E$ في $C_{3v}$ (واختزل الأخير).
\end{exercise}

\begin{exercise}[$\star$]\label{exo:b3:group-theory-applied:4}
% ledger: vib:CH4.nu1, vib:CH4.nu2, vib:CH4.nu3, vib:CH4.nu4
للميثان \ce{CH4} انتقالات أساسية عند 2917 ($a_1$) و1534 ($e$) و3019 و1306 (كلاهما $t_2$)
\unit{cm^{-1}}. أيها يظهر في طيف تحت الأحمر، وأيها في طيف
رامان؟
\end{exercise}

\begin{exercise}[$\star\star$]\label{exo:b3:group-theory-applied:5}
من أجل \ce{BF3}، استنتج $\Gamma_{\mathrm{vib}} = A_1' + 2E' + A_2''$ من الذرات الباقية في
مواضعها، ثم عُدّ حزم تحت الأحمر ورامان. أي نمط لا يُرى إلا في أحد
الطيفين، وأيها يُرى في كليهما؟
\end{exercise}

\begin{exercise}[$\star\star$]\label{exo:b3:group-theory-applied:6}
كم حزمة في تحت الأحمر وفي رامان يُظهر \ce{SF6}؟ وأي نمط صامت؟ ولماذا
لا تظهر أي حزمة في الطيفين معًا؟
\end{exercise}

\begin{exercise}[$\star\star$]\label{exo:b3:group-theory-applied:7}
طبّق \omterm{def:b3:group-theory-applied:projection}{مؤثر الإسقاط} على $h_2$ من أجل التمثيل $E$ للزمرة $C_{3v}$ 
وبيّن أن معامدة الناتج مع $2h_1 - h_2 - h_3$ تعطي $h_2 - h_3$.
\end{exercise}

\begin{exercise}[$\star\star$]\label{exo:b3:group-theory-applied:8}
تنبأ بعدد حزم CO في تحت الأحمر للمعقد \ce{ML3(CO)3} بصورتيه fac وmer،
وللمعقد \ce{M(CO)6}.
\end{exercise}

\begin{exercise}[$\star\star$]\label{exo:b3:group-theory-applied:9}
بيّن أن مدارات $\sigma$ الستة لربائط معقد ثماني الأوجه تولّد
$A_{1g} + E_g + T_{1u}$، وبيّن أي مدارات الفلز يمكن أن يتحد مع كل منها.
\end{exercise}

\begin{exercise}[$\star\star\star$]\label{exo:b3:group-theory-applied:10}
للتمثيل المنتظم $\chi(E) = h$ و$\chi(R) = 0$ من أجل $R \ne E$. بيّن
\omterm{thm:b3:group-theory-applied:reduction}{بصيغة الاختزال} أنه يحتوي كل \omterm{def:b3:point-groups:irrep}{تمثيل غير قابل للاختزال}
$\Gamma_i$ بالضبط $d_i$ مرة، واستنتج $\sum_id_i^2 = h$.
\end{exercise}

\begin{exercise}[$\star\star\star$]\label{exo:b3:group-theory-applied:11}
الأسيتيلين \ce{HC#CH} خطي. بالعمل في $D_{2h}$، جد $\Gamma_{\mathrm{vib}}$، وأعد تجميعه
برموز $D_{\infty h}$، وتنبأ بطيفي تحت الأحمر ورامان.
\end{exercise}

\begin{exercise}[$\star\star\star$]\label{exo:b3:group-theory-applied:12}
في الماء (المستوي $xz$)، أيّ ترقيات الإلكترون الواحد هذه تعطي انتقالات
مسموحة، وبأي استقطاب: $1b_2 \to 4a_1$، $3a_1 \to 4a_1$،
$1b_2 \to 2b_1$، $1b_1 \to 4a_1$؟
\end{exercise}

\section{مسألة: cis أم trans؟ معقد كربونيلي يُقرأ من طيفه}

\begin{problem}[{مسألة نهاية الأسبوع --- الزمر النقطية لأربعة متماكبات وتمثيلات
تمدد CO فيها، وحزمها في تحت الأحمر ورامان، وتحديد هوية
ناتج من طيفه المقيس}]
\label{pb:b3:group-theory-applied:1}
% ledger: ct:C2v, ct:C3v, ct:D4h
يحمل فلز M ثماني الأوجه ربائط كربونيلية وربائط أخرى L (تُعدّ كل L
نقطة). ونأخذ في الحسبان أربعة متماكبات: \ce{ML4(CO)2} بصورتيه cis
وtrans، و\ce{ML3(CO)3} بصورتيه fac وmer. ونفترض أن الربائط L لا
تخفض التناظر أكثر مما تفرضه مواضعها. يعطي تخليقٌ للمعقد
\ce{ML3(CO)3} ناتجًا يُظهر طيفه في تحت الأحمر ثلاث حزم قوية
عند 2048 و1965 و\qty{1938}{cm^{-1}} (معطيات المسألة)؛ ويُظهر طيفه في
رامان ثلاث حزم في المواضع نفسها تقريبًا.

\medskip\noindent\textbf{الجزء الأول --- الزمر النقطية.}
\begin{enumerate}
  \item ارسم المتماكبات الأربعة على ثماني أوجه.
  \item أعطِ \omterm{def:b3:point-groups:point-group}{الزمرة النقطية} للمعقد cis-\ce{ML4(CO)2}.
  \item وللمعقد trans-\ce{ML4(CO)2}.
  \item وللمعقد fac-\ce{ML3(CO)3}.
  \item وللمعقد mer-\ce{ML3(CO)3}.
  \item أي المتماكبات الأربعة له \omterm{def:b3:point-groups:inversion}{مركز انقلاب}؟
\end{enumerate}

\medskip\noindent\textbf{الجزء الثاني --- تمددات CO.}
\begin{enumerate}[resume]
  \item فسّر لماذا تحمل متجهات الروابط C--O تمثيلًا تساهم فيه
        العملية بالعدد 1 عن كل CO يبقى في موضعه.
  \item جد $\Gamma_{\mathrm{CO}}$ للمتماكب cis.
  \item وللمتماكب trans.
  \item وللمتماكب fac.
  \item وللمتماكب mer.
  \item قارن كل بعد بعدد ربائط CO.
\end{enumerate}

\medskip\noindent\textbf{الجزء الثالث --- النشاط.}
\begin{enumerate}[resume]
  \item عُدّ تمددات CO النشطة في تحت الأحمر لكل متماكب.
  \item عُدّ التمددات النشطة في رامان.
  \item أي متماكب تظهر فيه \omterm{prop:b3:group-theory-applied:mutual-exclusion}{قاعدة الاستبعاد المتبادل}؟
  \item في المتماكب trans، صِف النمط \omterm{def:b3:group-theory-applied:activity}{النشط في تحت الأحمر}: هل يتمدد
        الكربونيلان بالطور نفسه أم بطورين متعاكسين؟
  \item في المتماكب fac، لماذا لا تعطي مجموعات CO الثلاث إلا حزمتين؟
  \item ماذا يخبرك عدد حزم يساوي 1 (في تحت الأحمر)؟
\end{enumerate}

\medskip\noindent\textbf{الجزء الرابع --- الناتج.}
\begin{enumerate}[resume]
  \item أي متماكب يدل عليه طيف تحت الأحمر للناتج؟
  \item هل يتوافق طيف رامان معه؟
  \item الحزمة الأعلى هي التمدد المتطاور لمجموعات CO. إلى أي
        \omterm{def:b3:point-groups:irrep}{تمثيل غير قابل للاختزال} تنتمي في ذلك المتماكب؟
  \item هل يمكن أن يأتي الطيف من خليط من المتماكبين؟ وأي
        قياس إضافي يحسم الأمر؟
  \item اذكر النتيجة: عدد تمددات CO النشطة في تحت الأحمر للمتماكب mer،
        مقابل عددها للمتماكب fac.
\end{enumerate}
\end{problem}
